\subsection{長さ2・1・1の全必要被覆とH3排除}
\label{subsec:cubic-CY-finite211-cover}
\begin{theorem}\label{thm:cubic-CY-finite211-empty}
零交換子CY包絡を持つH3の対称射影の基点型は211にならない。
従って残る必要対称射影型は1111のみである。
\end{theorem}
\begin{proof}
固定出典\cite{AS4CYpinned}の三つの必要被覆、すなわち
非零第一層(DS1)--(DS10)、非零第二層(DT4)--(DT11)、最終零層(DZ1)--(DZ20)を用いる。
全AS正則性の十分性や全自己同型群は輸入しない。
元の対称射影と印付けは
\[
 U=(b^2,ac,ad,bc,bd,cd),\qquad
 Q=-z_1z_4+z_2z_3-z_0z_5+z_3z_4.
\]
全21混合係数の関係 $R=U+KB^{\mathsf T}A^{-1}$ と
Schur条件 $D=C+gJ-B^{\mathsf T}A^{-1}B=0$ を扱う。
基点環 $k[t]/t^2\times k\times k$ の全切断への作用から、全印付き群は
\[
 N(p)=I+pE_{10},\quad T(v,w)=\operatorname{diag}(1,1,v,w),\quad
 W=\begin{pmatrix}1&0&0&0\\0&-1&0&0\\0&0&0&1\\0&0&1&0\end{pmatrix}
\]
と共通スカラーで生成される。
$D_{00}=4x_0x_{15}$ なので、以下の三層が全てを覆う。

\paragraph{第一の非零層 $x_0\ne0$}
$x_0=1,x_{15}=0$ とする。
$(x_3,x_5)\ne(0,0)$ では交換と尺度変更で $x_3=1$、
$N$ の単位pivotで $x_7=0$ とできる。
残余式 $4x_4(x_4^2-x_5)=0$ から、$x_4\ne0$ は
$\mathcal L_{\pm i/4}$、$x_4=0,x_5\ne0$ は $\mathcal T_h$、
両方零は空となる。$x_3=x_5=0,x_4\ne0$ も空idealである。
従って残る上部ブロックは零。
$(x_7,x_{12})\ne(0,0)$ では $x_7=1,x_8=0$ とし、
$u=x_{11}=x_{13},s=x_{12},r=x_9,z=x_{20}$ に対する全残余式
\[
 z=s-u^2+1/16,\quad x_1z=x_2z=0,\quad (u-r)z=0,\quad rz=(u+r)/16
\]
を得る。$z=0$ は実平方 $b^2$、それ以外は $\mathcal U_u$ 又は $\mathcal V_h$。
$x_7=x_{12}=0$ の枝は、平方を避けると $\mathcal P$ だけである。
これらの全表示は既存のDS211節の全パラメータ排除で消える。

\paragraph{第二の非零層 $x_0=0,x_{15}\ne0$}
$t=x_{15}$ を単元とする。
$(x_6,x_{10})\ne(0,0)$ の三必要chartは全て空ideal。
上部ブロック非零の残余は二つの二変数族で三次商次元21となり排除される。
零ブロックの場合、$l=x_{11}-t,r=x_9,d=x_{13}-x_9$ とすると
\[
 ld-l^2=1/16,\qquad (16l^2-1)(2l+t)=0.
\]
$l=\pm1/4$ は、一般枝と二つの境界族が全て三次商次元22で排除される。
残りは $t=-2l,l\ne0,\pm1/4$。
全混合残余式から $x_{16}=x_{17}=x_{18}=x_{19}=0$ と
$x_{20}(16l^2+1)=0$ が従う。
従って既存の二パラメータ表示 $\mathcal A_{l,r}$ 又は例外表示 $\mathcal B_r$ に尽くされ、
定理\ref{thm:ds211-A-all-excluded}と定理\ref{thm:ds211-B-all-excluded}により全排除される。

\paragraph{最終零層 $x_0=x_{15}=0$}
この層の全分岐には高次障害が必要である。
$(x_6,x_{10})\ne(0,0)$ なら、全群で $x_6=1,x_1=0$ とする。
$x_{10}=0$ は空。$x_{10}=1$ では上部三係数が $\kappa=\pm i/8$ に等しい。
$x_7=0$ と $x_7\ne0$ の全因数分岐はそれぞれ一般一変数族と二つの代数的点に尽くされる。
一般族の五次商次元は57、四点の四次商次元は36である。
共役 $\kappa$ にも同じrank障害が成立する。

以後 $x_6=x_{10}=0$。
上部ブロックで $x_4\ne0$ は空ideal、$x_4=0$ で対角非零の場合は
二つの一変数族 $M^0_{h,\sigma},M^1_{h,\sigma}$ に移る。
$h\ne0,\sigma=\pm i/2$ に対し
\[
 (\dim C_0,\ldots,\dim C_7)=(1,4,10,20,35,56,84,123)
\]
なので、七次の必要次元120に反する。六次までの一致を正則性の証拠とはしない。

上部ブロックも零なら、全残余式から
$x_{11}^2-x_7x_{12}=1/16$ が従う。
$x_7=x_{12}=0$ の五必要代表・族は三次商次元22,23,23,21,21。
それ以外は $x_7=1,x_8=0$ とし、$(x_1,x_2)=(0,0),(0,1),(1,x_2)$ に分ける。
最初の場合の非平方枝は四次39、三次23の二点、五次60の共役点。
第二の場合は全因数分岐が四次39の二族と四次42の一族。
最後の場合は $z=x_{20}\ne0$ 又は $z=0,u=x_{18}\ne0$ の二つの二変数族で、
四次商次元はそれぞれ39と42である。
前者の元分母は $z$ と $u^2-z$、後者は二つの明記した非零係数の冪だけである。
$z=u^2\ne0$ は残余式自体に矛盾し、$z=u=0$ は実平方である。
これで最終零層の全分岐が尽くされる。

全除外族では元関係の分母を明記した単元だけに限定する。
元の全語関係行列のrankは閉値への特殊化で増えないので、
分数体上の商次元下界が全許容値に及ぶ。
当repoの独立商計算は最終零層22代表・族と第二層5族を元六関係から検証する。
残る第一層の表示は
定理\ref{thm:ds211-TU-fixed-excluded}と定理\ref{thm:ds211-V-all-excluded}の表示と一致する。
全てH3でないので結論を得る。
\end{proof}

公開再現器は\\
\texttt{verify\_cubic\_CY\_finite211\_cover.py}。
全分岐の必要被覆は固定出典の通常証明と上流全control再実行に依存し、
当repoの独立商計算、AS十分性、相対分類と区別する。
